% Source PDF pages 34--35; manual pages 2-5--2-6.
\subsection{Earth-Centered and Fixed, Cartesian (ECFC) Coordinate System}\label{sec:ecfc}

An earth-centered and fixed, Cartesian coordinate system is used for the equations
of motion and is therefore the primary coordinate system in TAOS. Its origin is at
the center of the earth. The $x$ and $y$ axes are in the equatorial plane, and the
$z$ axis points through the north pole, as shown in
\cref{fig:ecfc-coordinate-system}. The earth spins about the $z$ axis. The $x$ axis
is aligned with the Greenwich meridian ($\lambda=0^{\circ}$), and the $y$ axis is
aligned with the $90^{\circ}$ east meridian. Unit vectors
$\hat{x}_{\oplus}$, $\hat{y}_{\oplus}$, and $\hat{z}_{\oplus}$ are parallel to
these axes.

\begin{figure}[htbp]
  \centering
  \resizebox{0.58\linewidth}{!}{\begin{tikzpicture}[scale=0.95,>=Latex,line cap=round,line join=round]
  \coordinate (O) at (0,0);
  \begin{scope}
    \clip (O) circle (2.05);
    % Latitude grid.
    \draw[thin] (-2.00,0.00) arc[start angle=180,end angle=360,x radius=2.00,y radius=0.36];
    \draw[thin] (-2.00,0.00) arc[start angle=180,end angle=0,x radius=2.00,y radius=0.36];
    \draw[thin] (-1.88,0.78) arc[start angle=170,end angle=370,x radius=1.90,y radius=0.47];
    \draw[thin] (-1.88,-0.78) arc[start angle=190,end angle=350,x radius=1.90,y radius=0.47];
    \draw[thin] (-1.52,1.38) arc[start angle=158,end angle=382,x radius=1.58,y radius=0.34];
    \draw[thin] (-1.52,-1.38) arc[start angle=202,end angle=338,x radius=1.58,y radius=0.34];
    % Meridians.
    \draw[thin] (0,2.05) arc[start angle=90,end angle=270,x radius=0.52,y radius=2.05];
    \draw[thin] (0,2.05) arc[start angle=90,end angle=-90,x radius=0.52,y radius=2.05];
    \draw[thin] (0,2.05) arc[start angle=90,end angle=270,x radius=1.10,y radius=2.05];
    \draw[thin] (0,2.05) arc[start angle=90,end angle=-90,x radius=1.10,y radius=2.05];
    \draw[thin] (0,2.05) arc[start angle=90,end angle=270,x radius=1.62,y radius=2.05];
    \draw[thin] (0,2.05) arc[start angle=90,end angle=-90,x radius=1.62,y radius=2.05];

    % Simplified continental outlines. These are intentionally schematic but retain
    % the geographic cue present in the source illustration.
    \draw[line width=0.48pt,fill=black!9]
      (-0.53,1.09) .. controls (-0.18,1.30) and (0.18,1.20) .. (0.45,0.96)
      .. controls (0.72,0.74) and (0.72,0.48) .. (0.52,0.28)
      .. controls (0.35,0.10) and (0.30,-0.05) .. (0.22,-0.34)
      .. controls (0.12,-0.73) and (-0.11,-1.05) .. (-0.35,-1.12)
      .. controls (-0.50,-0.92) and (-0.67,-0.68) .. (-0.71,-0.38)
      .. controls (-0.76,-0.02) and (-0.69,0.31) .. (-0.78,0.54)
      .. controls (-0.73,0.79) and (-0.64,0.94) .. (-0.53,1.09);
    \draw[line width=0.48pt,fill=black!9]
      (-0.60,1.16) .. controls (-0.89,1.26) and (-1.05,1.48) .. (-1.29,1.45)
      .. controls (-1.50,1.39) and (-1.55,1.18) .. (-1.37,1.02)
      .. controls (-1.18,0.91) and (-0.99,0.97) .. (-0.83,0.87);
    \draw[line width=0.48pt,fill=black!9]
      (0.37,1.03) .. controls (0.72,1.31) and (1.14,1.38) .. (1.54,1.18)
      .. controls (1.78,1.02) and (1.88,0.72) .. (1.63,0.55)
      .. controls (1.38,0.39) and (1.11,0.49) .. (0.91,0.30)
      .. controls (0.78,0.17) and (0.84,-0.04) .. (0.67,-0.15);
    \draw[line width=0.48pt,fill=black!9]
      (-1.42,0.62) .. controls (-1.72,0.80) and (-1.89,0.58) .. (-1.72,0.34)
      .. controls (-1.56,0.15) and (-1.31,0.12) .. (-1.23,-0.11)
      .. controls (-1.17,-0.31) and (-1.38,-0.43) .. (-1.27,-0.63)
      .. controls (-1.15,-0.86) and (-0.95,-1.08) .. (-0.82,-1.37)
      .. controls (-0.63,-1.24) and (-0.60,-0.98) .. (-0.71,-0.76)
      .. controls (-0.82,-0.52) and (-0.73,-0.29) .. (-0.82,-0.09)
      .. controls (-0.93,0.17) and (-1.15,0.28) .. (-1.42,0.62);
    \draw[line width=0.48pt,fill=black!9]
      (1.07,-1.04) .. controls (1.27,-0.91) and (1.56,-0.94) .. (1.66,-1.15)
      .. controls (1.54,-1.34) and (1.27,-1.42) .. (1.05,-1.31)
      .. controls (0.93,-1.22) and (0.95,-1.10) .. (1.07,-1.04);
  \end{scope}

  \draw[line width=0.95pt] (O) circle (2.05);
  \draw[->,line width=1.05pt] (O) -- (-2.72,-1.16) node[below left] {$\hat{x}_{\oplus}$};
  \draw[->,line width=1.05pt] (O) -- (2.92,0.00) node[right] {$\hat{y}_{\oplus}$};
  \draw[->,line width=1.05pt] (O) -- (0,2.82) node[above] {$\hat{z}_{\oplus}$};
  \fill (O) circle (1.35pt);
\end{tikzpicture}
}
  \caption{The ECFC Coordinate System.}
  \label{fig:ecfc-coordinate-system}
\end{figure}

The following output variables give position, velocity, and acceleration components
in ECFC coordinates.

\taossubhead{Position}
\noindent\begin{tabularx}{\textwidth}{@{}>{\ttfamily}l>{$\displaystyle}l<{$}X@{}}
  xecfc & \vec{r}\cdot\hat{x}_{\oplus} &
    Component of the position vector in the ECFC $x$ direction. The position
    vector extends from the earth's center to the vehicle's center of mass. \\
  yecfc & \vec{r}\cdot\hat{y}_{\oplus} &
    Component of the position vector in the ECFC $y$ direction. \\
  zecfc & \vec{r}\cdot\hat{z}_{\oplus} &
    Component of the position vector in the ECFC $z$ direction. \\
\end{tabularx}

\taossubhead{Velocity}
\noindent\begin{tabularx}{\textwidth}{@{}>{\ttfamily}l>{$\displaystyle}l<{$}X@{}}
  xecfcdt & \vec{V}_{\oplus}\cdot\hat{x}_{\oplus} &
    Component of earth-relative velocity in the ECFC $x$ direction. \\
  yecfcdt & \vec{V}_{\oplus}\cdot\hat{y}_{\oplus} &
    Component of earth-relative velocity in the ECFC $y$ direction. \\
  zecfcdt & \vec{V}_{\oplus}\cdot\hat{z}_{\oplus} &
    Component of earth-relative velocity in the ECFC $z$ direction. \\
\end{tabularx}

\taossubhead{Acceleration}
\noindent\begin{tabularx}{\textwidth}{@{}>{\ttfamily}l>{$\displaystyle}l<{$}X@{}}
  xecfcdt2 & \vec{a}_{\oplus}\cdot\hat{x}_{\oplus} &
    Component of earth-relative acceleration in the ECFC $x$ direction. \\
  yecfcdt2 & \vec{a}_{\oplus}\cdot\hat{y}_{\oplus} &
    Component of earth-relative acceleration in the ECFC $y$ direction. \\
  zecfcdt2 & \vec{a}_{\oplus}\cdot\hat{z}_{\oplus} &
    Component of earth-relative acceleration in the ECFC $z$ direction. \\
\end{tabularx}

Because the equations of motion are written in ECFC coordinates, these quantities
are obtained directly from numerical integration, as described in
\cref{sec:numerical-integration}.
